\documentclass[10pt,a4paper]{amsart}
\usepackage[margin=19mm]{geometry}
\usepackage{amsmath,amssymb,mathtools}
\usepackage[scr=rsfs,cal=euler]{mathalfa}
\usepackage{xfrac}
\usepackage[unicode,hidelinks,bookmarks=false]{hyperref}
\newcommand{\ie}{i.e.\ }
\newcommand{\cf}{cf.\ }
\newcommand{\resp}{respectively\ }
\newcommand{\Subsection}{\subsection}
\providecommand{\nobreakdash}{\nobreak\hbox{-}}
\setlength{\emergencystretch}{2em}
\multlinegap0pt
\usepackage{amssymb}
\usepackage{enumerate}
\usepackage{tikz}
\usepackage{xcolor}
\newtheorem{thm}{Theorem}
\newtheorem{lem}{Lemma}
\newcommand{\R}{\mathbf{R}}
\newcommand{\e}{\varepsilon}
\title{Blow-up of solutions to the Euler-Poisson-Darbox equation with critical power nonlinearity}
\author{Mengting Fan, Ning-An Lai, Hiroyuki Takamura}
\date{Source: arXiv:2603.19614v2 (CC BY 4.0)}
\begin{document}
\maketitle

\noindent\emph{Source and licence.} Blow-up of solutions to the Euler-Poisson-Darbox equation with critical power nonlinearity. Version arXiv:2603.19614v2 (CC BY 4.0), distributed by arXiv under the Creative Commons Attribution 4.0 International licence, http://creativecommons.org/licenses/by/4.0/, as recorded in the arXiv licence metadata for this exact version and shown on the abstract page https://arxiv.org/abs/2603.19614v2. Full licence notice: \texttt{licenses/arxiv-2603.19614-CC-BY-4.0.txt}.\par\smallskip

\noindent\textit{Editorial scope.} Counted unit: the article's lem lbq (source0.tex lines 572--588). The article's complete original proof of it is delivered with it (source0.tex lines 589--741, one \texttt{proof} environment of 153 lines). No mathematics is written, repaired or re-derived: the statement and the proof are the article's own lines, in the article's own notation, and no retained line is substituted or re-flowed.  Retained uncounted as the source-local context the proof needs: lines 126--134; lines 231--265; lines 502--508; lines 530--535; lines 546--550; lines 560--564.  Nothing was omitted from any retained span, so the package carries no omission record.  No retained line contains a citation, so the wrapper carries no bibliography and no bibliography entry.  No retained line contains a figure environment, an image-inclusion command or drawing code, so no image is delivered and the package declares no asset dependency.  Editorial formatting only, in addition: an AMS \texttt{amsart} preamble that restates the article's own declarations and macros used by the retained lines, namely the environments \texttt{thm}, \texttt{lem} on separate counters, together with the standard packages those lines need.  The retained environments print in this document as Theorem 1, Lemma 1 in order of appearance, whereas the article prints the same items with the numbering of its own full document.  The complete native source of the article as distributed in the arXiv e-print for this exact version is bundled unchanged as \texttt{sources/196788/source0.tex}.
\begin{equation}
\label{SHEPD}
\left\{
\begin{aligned}
& u_{tt}-\Delta u+\frac{\mu}{t}u_t=t^\alpha|u|^p, ~~~~t>0, x\in \R^n,\\
& u(0, x)=u_0(x), \quad u_t(0, x)=0,~~~~x\in \R^n
\end{aligned}
\right.
\end{equation}

\begin{thm}
\label{hddamped}
Let
\begin{equation}
\label{al}
\left\{
\begin{aligned}
& \alpha\ge 0,~~~n\ge 3,\\
& \alpha>0,~~~~n=2
\end{aligned}
\right.
\end{equation}
and
\[
0<\mu\leq\mu^*(n,\alpha), p=p_S(n+\mu, \alpha),
\]
 where
\begin{equation}
\begin{aligned}
\mu^*(n,\alpha)=\frac{2n^2+(n+2+\alpha)(n\alpha+2+\alpha)}{(n+2+\alpha)(2+\alpha)}.
\end{aligned}
\end{equation}
Assume the initial data $u_0(x)\in H^1(\R^n)$ is positive and has compact support
\begin{equation}
\begin{aligned}
supp~u_0(x)\subset\{|x|\le 1\},
\end{aligned}
\end{equation}
then there exists a constant $\e_0=\e_0(u_0,n,p,\mu,\alpha)>0$ such that the lifespan of the energy solution to the Cauchy problem \eqref{SHEPD} satisfies
\begin{equation}
\label{lf2}
T(\e)\le \exp\left(C\e^{-p(p-1)}\right),
\end{equation}
where $C$ denotes a generic positive constant independent of $\e$ which may have different values from line to line.
\end{thm}

\begin{equation}\label{beode}
\begin{aligned}
& \lim_{t\rightarrow 0}\left(-h'(t)+\mu\frac{h(t)}{t}\right)\thicksim 2^{\frac{\mu-1}{2}}\Gamma\left(\frac{\mu+1}{2}\right)\triangleq C_0>0,\\
&h(t)\thicksim t^{\frac{\mu}2}e^{-t}\quad\mbox{for}\ t>R,\\
&|h'(t)|\thicksim t^{\frac{\mu}2}e^{-t}\quad\mbox{for}\ t>R,\\
\end{aligned}
\end{equation}

\begin{equation}\label{rela}
\begin{aligned}
&K_\nu'(z)=-\frac12\left[K_{\nu-1}(z)+K_{\nu+1}(z)\right],\\
&K_\nu(z)=-\frac{z}{2\nu}\left[K_{\nu-1}(z)-K_{\nu+1}(z)\right].\\
\end{aligned}
\end{equation}

\begin{equation}\label{testpro}
\begin{aligned}
0<\phi(x)\le C(1+|x|)^{-\frac{n-1}{2}}e^{|x|}.
\end{aligned}
\end{equation}

\begin{equation}\label{bq}
\begin{aligned}
b_q(t,x)=\int_{0}^{1}h_\lambda(t)\phi_{\lambda}(x)\lambda^{q-1}d\lambda,
\end{aligned}
\end{equation}

\begin{lem}\label{lbq}
The function $b_q(t, x)$ defined in \eqref{bq} satisfies
\begin{equation}\label{eqforbq}
\partial_t^2b_q-\Delta b_q-\partial_t\left(\frac{\mu}{t}b_q\right)=0,~~~t>0, x\in\R^n,
\end{equation}
\begin{equation}\label{asymforbq}
\begin{aligned}
&b_q(t, x)\thicksim t^{-q},~~~t\ge 1,
\end{aligned}
\end{equation}
and
\begin{equation}\label{asymfortbq}
\begin{aligned}
&\partial_t b_q(t, x)\lesssim t^{\frac{\mu}{2}-\frac{n-1}{2}}(t+2-r)^{\frac{1}{p}-1},~~~t\ge 1.
\end{aligned}
\end{equation}
\end{lem}

\begin{proof}
We may compute $h'(\lambda t)$ by using \eqref{rela}
\begin{equation}\label{dh}
\begin{aligned}
&h'(\lambda t)\\
=&\frac{\mu+1}{2}(\lambda t)^{\frac{\mu-1}{2}}\lambda K_{\frac{\mu-1}{2}}(\lambda t)
+(\lambda t)^{\frac{\mu+1}{2}}\lambda K'_{\frac{\mu-1}{2}}(\lambda t)\\
=&\frac{\mu+1}{2}(\lambda t)^{\frac{\mu-1}{2}}\lambda K_{\frac{\mu-1}{2}}(\lambda t)+(\lambda t)^{\frac{\mu+1}{2}}\lambda
\left(-\frac12K_{\frac{\mu-3}{2}}(\lambda t)-\frac12K_{\frac{\mu+1}{2}}(\lambda t)\right)\\
=&\frac{\mu+1}{2}(\lambda t)^{\frac{\mu-1}{2}}\lambda K_{\frac{\mu-1}{2}}(\lambda t)-\frac12(\lambda t)^{\frac{\mu+1}{2}}\lambda
\left(K_{\frac{\mu+1}{2}}(\lambda t)-\frac{2\cdot\frac{\mu-1}{2}}{\lambda t}
K_{\frac{\mu-1}{2}}(\lambda t)\right)\\
&-\frac12(\lambda t)^{\frac{\mu+1}{2}}\lambda K_{\frac{\mu+1}{2}}(\lambda t)\\
=&\mu (\lambda t)^{\frac{\mu-1}{2}}\lambda K_{\frac{\mu-1}{2}}(\lambda t)-(\lambda t)^{\frac{\mu+1}{2}}\lambda K_{\frac{\mu+1}{2}}(\lambda t).\\
\end{aligned}
\end{equation}
From the definition of $h_\lambda(t)$ and $b_q(t, x)$, direct computation yields
\begin{equation}\label{l41}
\begin{aligned}
\partial_tb_q(t,x)&=\int_{0}^{1}h_\lambda'(t)\phi_\lambda(x)\lambda^{q-1}d\lambda\\
&=\frac{\mu}{t}\int_{0}^{1}(\lambda t)^{\frac{\mu+1}{2}}K_{\frac{\mu-1}{2}}(\lambda t)\phi_\lambda(x)\lambda^{q-1}d\lambda\\
&-\int_{0}^{1}(\lambda t)^{\frac{\mu+1}{2}}K_{\frac{\mu+1}{2}}(\lambda t)\phi_\lambda(x)\lambda^{q}d\lambda.
\end{aligned}
\end{equation}
In a similar way, we have
\begin{equation}\label{l42}
\begin{aligned}
\partial_t^2b_q=&-\frac{\mu}{t^2}b_q+\frac{\mu}t\int_0^1\partial_th_\lambda(t)
\phi_\lambda(x)\lambda^{q-1}d\lambda\\
&-\frac{\mu+1}{2}\int_0^1(\lambda t)^{\frac{\mu-1}{2}}K_{\frac{\mu+1}{2}}(\lambda t)\phi_\lambda(x)\lambda^{q+1}d\lambda\\
&-\int_0^1(\lambda t)^{\frac{\mu+1}{2}}\partial_tK_{\frac{\mu+1}{2}}(\lambda t)\phi_\lambda(x)\lambda^{q}d\lambda\\
=&-\frac{\mu}{t^2}b_q+\frac{\mu^2}{t^2}b_q-\frac{\mu}t\int_0^1(\lambda t)^{\frac{\mu+1}{2}}K_{\frac{\mu+1}{2}}(\lambda t)\phi_\lambda(x)\lambda^{q}d\lambda\\
&+\int_0^1(\lambda t)^{\frac{\mu+1}{2}}K_{\frac{\mu-1}{2}}(\lambda t)\phi_\lambda(x)\lambda^{q+1}d\lambda.\\
\end{aligned}
\end{equation}
Noting that
\begin{equation}\label{l42a}
\begin{aligned}
\Delta b_q=\int_0^1(\lambda t)^{\frac{\mu+1}{2}}K_{\frac{\mu-1}{2}}(\lambda t)\phi_\lambda(x)\lambda^{q+1}d\lambda,
\end{aligned}
\end{equation}
then \eqref{eqforbq} follows by combining \eqref{l41}, \eqref{l42} and \eqref{l42a}.

We next show the asymptotic behavior \eqref{asymforbq}. For $t\ge 1$, by $\eqref{beode}_2$, one has
\begin{equation}\label{l43}
\begin{aligned}
b_q(t, x)\gtrsim& \int_{\frac{1}{2(t+1)}}^{\frac{1}{t+1}}(\lambda t)^{\frac{\mu}{2}}
e^{-\lambda t}\lambda^{q-1}d\lambda\\
\gtrsim& t^{\frac{\mu}2}\int_{\frac{1}{2(t+1)}}^{\frac{1}{t+1}}\lambda^{q-1+\frac{\mu}{2}}
e^{-\lambda (t+R)}d\lambda\\
\gtrsim& t^{-q}\int_{\frac12}^1s^{q+\frac {\mu}{2}-1}e^{-s}ds\\
\gtrsim& t^{-q}.\\
\end{aligned}
\end{equation}
In the opposite direction, we divide the proof into two parts. For $r\le \frac{t+1}{2}$, it is easy to get by $\eqref{beode}_2$ and \eqref{testpro}
\begin{equation}\label{l44}
\begin{aligned}
b_q(t, x)\lesssim& \int_0^1(\lambda t)^{\frac{\mu}2}e^{-\lambda t}(1+\lambda r)^{-\frac{n-1}2}e^{\lambda r}\lambda^{q-1}d\lambda\\
\lesssim&t^{\frac{\mu}2}\int_0^1e^{-\frac{\lambda(t+1)}{2}}
\lambda^{q+\frac{\mu}2-1}d\lambda\\
\lesssim&t^{-q}\int_0^{\infty}e^{-s}s^{q+\frac{\mu}{2}-1}ds\\
\lesssim&t^{-q},
\end{aligned}
\end{equation}
while for $\frac{t+1}{2}\le r\le t+1$, it holds that
\begin{equation}\label{bigr}
\begin{aligned}
b_q(t, x)\lesssim& \int_0^1(\lambda t)^{\frac{\mu}2}e^{-\lambda t}(1+\lambda r)^{-\frac{n-1}2}e^{\lambda r}\lambda^{q-1}d\lambda\\
\lesssim&\int_0^1(\lambda t)^{\frac \mu 2}\bigg(1+\lambda(t+1)\bigg)^{-\frac{n-1}{2}}\lambda^{q-1} d\lambda\\
\lesssim&t^{-q}\int_0^{\infty}(1+s)^{-\frac{n-1}{2}}s^{\frac \mu 2+q-1}ds\\
\lesssim&t^{-q}\left(\int_0^{1}(1+s)^{-\frac{n-1}{2}}s^{\frac{n-1}2-\frac 1p-1}ds+\int_1^{\infty}(1+s)^{-\frac 1p-1}ds\right),\\
\end{aligned}
\end{equation}
where we use the fact that
\[
q=\frac{n-\mu-1}{2}-\frac{1}{p}.
\]
Obviously, the last term in \eqref{bigr} in integrable, while for the second to the last term, we should require $\frac{n-1}{2}>\frac 1p$ to ensure the integrability.
Actually it is equivalent to $p>\frac 2{n-1}, n\ge 2$, which is also equivalent to
\begin{equation}\label{eq1}
\begin{aligned}
\gamma\left(n, \mu, \alpha, \frac 2{n-1}\right)=4(n-1)^2+2(n-3)\mu+4(n-1)\alpha
>0,
\end{aligned}
\end{equation}
which always holds for $n\ge 3, \alpha\ge 0, \mu>0$. If $n=2$, \eqref{eq1} becomes to
\begin{equation}\label{eq2}
\begin{aligned}
\mu<2+2\alpha,
\end{aligned}
\end{equation}
which also holds for $0<\mu\le \mu^*(2, \alpha)$ and $\alpha>0$. And hence under the assumption of Theorem \ref{hddamped}, the inequality \eqref{bigr} becomes
\begin{equation}\label{bigr6}
\begin{aligned}
b_q(t, x)\lesssim t^{-q}.
\end{aligned}
\end{equation}
Then \eqref{asymforbq} follows by combining \eqref{l43}, \eqref{l44} and \eqref{bigr6}.

We finally show the asymptotic behavior \eqref{asymfortbq}. We can write \eqref{l41} as
\begin{equation}\label{l45}
\begin{aligned}
\partial_tb_q(t,x)&=I+II,
\end{aligned}
\end{equation}
with
\[
\begin{aligned}
I&:=\frac{\mu}{t}\int_{0}^{1}(\lambda t)^{\frac{\mu+1}{2}}K_{\frac{\mu-1}{2}}(\lambda t)\phi_\lambda(x)\lambda^{q-1}d\lambda,\\
II&:=-\int_{0}^{1}(\lambda t)^{\frac{\mu+1}{2}}K_{\frac{\mu+1}{2}}(\lambda t)\phi_\lambda(x)\lambda^{q}d\lambda.
\end{aligned}
\]
By combining $\eqref{beode}_2, \eqref{beode}_3$, \eqref{testpro}, it is easy to estimate $I$ as
\begin{equation}
\begin{aligned}
|I|\leq \frac{\mu}{t}\int_{0}^{1}h_\lambda(t)\phi_\lambda(x)\lambda^{q-1}d\lambda\lesssim t^{-q-1}.
\end{aligned}
\end{equation}
For $II$, we split the proof into two parts: $0<r\leq\frac{t+1}{2}$ and $\frac{t+1}{2}\leq r \leq t+1$. For the former case, we have
\begin{equation}
\begin{aligned}
|II|&\lesssim \int_{0}^{1}(\lambda t)^{\frac{\mu+1}{2}}(\lambda t)^{-\frac{1}{2}}e^{-\lambda t}(1+\lambda r)^{-\frac{n-1}{2}}e^{\lambda r}\lambda^{q}d\lambda\\
&\lesssim \int_{0}^{1}(\lambda t)^{\frac{\mu}{2}}e^{-\lambda(t-r)}(1+\lambda r)^{-\frac{n-1}{2}}\lambda^{q}d\lambda\\
&\lesssim t^{\frac{\mu}{2}}\int_{0}^{1}e^{-\lambda(t-r)}(1+\lambda r)^{-\frac{n-1}{2}}\lambda^{q+\frac{\mu}{2}}d\lambda\\
&\lesssim t^{\frac{\mu}{2}}\int_{0}^{1}e^{-\frac{\lambda(t+1)}{2}}\lambda^{q+\frac{\mu}{2}}d\lambda\\
&\lesssim t^{-q-1}\int_{0}^{\infty}e^{-s}s^{\frac{n-1}{2}-\frac{1}{p}}ds\\
&\lesssim t^{-\frac{n-\mu-1}{2}+\frac{1}{p}-1},
\end{aligned}
\end{equation}
where
\[
q=\frac{n-\mu-1}{2}-\frac{1}{p}.
\]
If $\frac{t+1}{2}\leq r \leq t+1$, we get
\begin{equation}
\begin{aligned}
|II|&\lesssim t^{\frac{\mu}{2}}\int_{0}^{1}e^{-\lambda(t+2-r)}\big(\lambda (t+1)\big)^{-\frac{n-1}{2}}\lambda^{q+\frac{\mu}{2}}d\lambda\\
&\lesssim t^{\frac{\mu}{2}}\int_{0}^{t+2-r}e^{-s}(t+1)^{-\frac{n-1}{2}}s^{-\frac{n-1}{2}}(t+2-r)^{\frac{n-1}{2}}s^{q+\frac{\mu}{2}}\\
&\quad\cdot(t+2-r)^{-q-\frac{\mu}{2}}(t+2-r)^{-1}ds\\
&\lesssim t^{\frac{\mu}{2}-\frac{n-1}{2}}(t+2-r)^{\frac{n-1}{2}-q-\frac{\mu}{2}-1}\int_{0}^{\infty}e^{-s}s^{-\frac{n-1}{2}+q+\frac{\mu}{2}}ds\\
&\lesssim t^{\frac{\mu}{2}-\frac{n-1}{2}}(t+2-r)^{\frac{1}{p}-1}\int_{0}^{\infty}e^{-s}s^{-\frac{1}{p}}ds\\
&\lesssim t^{\frac{\mu}{2}-\frac{n-1}{2}}(t+2-r)^{\frac{1}{p}-1}\bigg(\int_{0}^{1}e^{-s}s^{-\frac{1}{p}}ds+\int_{1}^{\infty}e^{-s}s^{-\frac{1}{p}}ds\bigg)\\
&\lesssim t^{\frac{\mu}{2}-\frac{n-1}{2}}(t+2-r)^{\frac{1}{p}-1}.
\end{aligned}
\end{equation}
Thus, we obtain
\begin{equation}
\begin{aligned}
|\partial_tb_q|\lesssim |I|+|II|\lesssim t^{\frac{\mu}{2}-\frac{n-1}{2}}(t+2-r)^{\frac{1}{p}-1},
\end{aligned}
\end{equation}
which leads to the asymptotic behavior \eqref{asymfortbq}.
\end{proof}

\end{document}
